% !TEX TS-program = pdflatex
% !TeX program = pdflatex
% !TEX encoding = UTF-8
% !TeX spellcheck = en_US

\documentclass[11pt, a4paper]{article}
%\usepackage{fullpage}
\usepackage[left=1cm,right=1cm,top=1cm,bottom=2cm]{geometry}
\usepackage[fleqn]{amsmath}
\usepackage{amssymb}
%\usepackage{indentfirst}
\usepackage[T1]{fontenc}
\usepackage[utf8]{inputenc}
\usepackage[french,english]{babel}
\usepackage{txfonts} 
\usepackage[]{graphicx}
\usepackage{multirow}
\usepackage{hyperref}
\usepackage{parskip}
\usepackage{multicol}
\usepackage{wrapfig}

\usepackage{turnstile}%Induction symbole

\usepackage{tikz}
\usetikzlibrary{arrows, automata}
\usetikzlibrary{decorations.pathmorphing}

\renewcommand{\baselinestretch}{1}

\setlength{\parindent}{24pt}


\begin{document}

%\selectlanguage {french}
%\pagestyle{empty} 

\noindent
\begin{tabular}{ll}
\multirow{3}{*}{\includegraphics[width=1.5cm]{../../../extra/logo/esi.nlp.pdf}} & \'Ecole national Supérieure d'Informatique\\
& 2\textsuperscript{nd} year second cycle (2CSSID)\\
& NLP: Natural Language Processing (2022-2023)
\end{tabular}\\[.25cm]
\noindent\rule{\textwidth}{2pt}\\[-0.5cm]
\begin{center}
{\LARGE \textbf{Lab03: Parsing}}
\begin{flushright}
	Abdelkrime Aries
\end{flushright}
\end{center}\vspace{-0.5cm}
\noindent\rule{\textwidth}{2pt}

We want to implement a simple parser both for lexical analysis and syntactic parsing. 
Our parser is based on CKY algorithm.
In CKY, the grammar's rules can be devised into two parts: generation rules and lexicon.
There latter are the rules which generates the words.
In this case, we will replace those with "part-of-speech tagging".

\section*{1. Program description}

Our objective is to use PoS tagging as a lexicon function (find out for each word its PoS).
Then, the result will be fed into CKY algorithm to create a parse tree.
A first design is to implement these classes:
\begin{itemize}
	\item \textbf{MEMM}: This class implements the different methods used to score a language's sentence.
	\textbf{\slshape A part of this class must be implemented}.
	\item \textbf{CKY}: This class is used to parse a given sentence based on the CKY algorithm and the PoS tagger we already trained.
	\textbf{\slshape A part of this class must be implemented}.
	\item \textbf{Parser}: This class is used to train the two models in order to learn the whole pipeline.
	\textbf{\slshape This is already implemented}.
\end{itemize}
The other objective is to implement a program dependency-free; i.e. it does not depend on any external API.


\subsection*{1.1. MEMM class}

This is our PoS tagger based on MEMM.
The class has these methods:
\begin{itemize}
	\item encode: which encodes the past word, the past tag, the current word and the future word into a singular vector.
	You have to implement this method. 
	We defined a function outside the class in order to encode one word at a time.
	This function must also be implemented.
	\item Others are described in the source code.
\end{itemize}

\subsubsection*{encode\_word function}

This function is used to encode one single word.
Given a word, the features are the following (in this precise order):
\begin{itemize}
	\item is the first letter uppercase? (yes:1, no:0)
	\item does the word contains a period (.)? (yes:1, no:0)
	\item Given a list of prefixes (given in the code), for each one, did the word start with it. 
	It is a vector equals to the number of prefixes.
	\item Given a list of suffixes (given in the code), for each one, did the word ends with it. 
	It is a vector equals to the number of suffixes.
\end{itemize}

\subsubsection*{MEMM.encode}

The resulted code is based on these features:
\begin{itemize}
	\item is the past tag equals to "$<t>$"? (yes:1, no:0)
	\item onehot encoding of the tag 
	\item is the past word equals to "$<s>$"? (yes:1, no:0)
	\item vector encoding of past word using the past function
	\item vector encoding of current word using the past function
	\item is the future word equals to "$</s>$"? (yes:1, no:0)
	\item vector encoding of future word using the past function
\end{itemize}

Some useful functions are afforded such as onehot 


\subsection*{1.2. CKY class}

In this class, you have to implement the parse method based on CKY algorithm.
It must use the Part-or-Speech tagger called "self.lex".


\section*{2. Questions}

Answer these questions at the beginning of your code, as comments:
\begin{enumerate}
	\item PoS: We note that the vector representing our input is too sparse (many zeros) which causes a bigger model.
	Propose a solution.
	
	\item PoS: We note that the model is not really accurate.
	It is far from being acceptable.
	Propose a solution to enhance MEMM prediction.
	
	\item CKY: What is the benefit and the downfall of using PoS in CKY parsing instead of a hand-prepared lexicon?
	
	\item CKY: Can we change CKY to handle unitary productions (A --> B)? If yes, how? If no, why?
	
\end{enumerate}


\section*{3. Students' Grading}

\begin{itemize}
	\item Duration: 1.5h (Homework)
	\item Grade
	\begin{itemize}
		\item \textbf{encode word function} (4pts) = first element (0.5pts) + second (0.5pts) + prefixes (1pts) + suffixes (1pt) + order (1pt).
		\item \textbf{MEMM.encode method} (4.5pts) = past tag (1pt) + past word (1pt) + current word (0.5pts) + future word (1pt) + order (1pt).
		\item \textbf{CKY.parse method} (5.5pts) = lexicon integration (1.5pt) + CKY table creation (4pts).
		\item \textbf{questions grade} (4pts) = 1pt for each question.
		\item \textbf{In time grade} (2pts) = you have to return the homework at time. 
		Each half an hour is decreased by 0.25.
		After 4 hours, you will have a 0.
	\end{itemize}
\end{itemize}

\end{document}
